\begin{statementbox}
	\textbf{\textcolor{CStatement}{条件}}
	\begin{description}[style=nextline, leftmargin=2.8em, labelsep=0.5em, itemsep=0.35em]
		\item[\textbf{\textcolor{CStatement}{条件1（初值满足）：}}] $a_1 = f(1)$.
		\item[\textbf{\textcolor{CStatement}{条件2（递推保持）：}}] 对任意 $k\in\mathbb{N}^*$，若 $a_k = f(k)$ 成立，则能推出 $a_{k+1} = f(k+1)$ 也成立.
	\end{description}

	\textbf{\textcolor{CStatement}{结论}}
	\begin{description}[style=nextline, leftmargin=2.8em, labelsep=0.5em, itemsep=0.45em]
		\item[\textbf{\textcolor{CStatement}{结论一（通项恒等）：}}]
			\textbf{\textcolor{CStatement}{直观描述：}}
			从第一项开始数列的每一项都与猜想表达式一致.
			\[
				a_n = f(n)
				\qquad
				(n \in \mathbb{N}^*)
			\]

		\item[\textbf{\textcolor{CStatement}{推广结论（一般归纳原理）：}}]
			\textbf{\textcolor{CStatement}{直观描述：}}
			数学归纳法可用于证明任何与自然数有关的命题，不限于数列通项.
			\[
				P(1) \land (P(k) \Rightarrow P(k+1)) \Rightarrow \forall n\in\mathbb{N}^*,\ P(n)
			\]
	\end{description}

	\textbf{\textcolor{CStatement}{使用提醒：}}
	必须先通过观察前几项或递推关系猜想出通项公式 $f(n)$，否则本结论无法直接使用.
\end{statementbox}
